\section{Related Work}~\label{related}
Our work connects three lines of research: structured models for robot decision making, planning under uncertainty, and interactive acquisition of external information. We organize this review around how uncertainty is represented, when additional information is requested, and what the user is asked to provide.

\subsection{Structured Planning under Symbolic State Uncertainty}~\label{rel:structured_uncertainty}
Knowledge-based robotic systems represent task-relevant objects, attributes, relations, goals, and constraints using symbolic predicates~\cite{fikes1971strips,aeronautiques1998pddl}, ontologies~\cite{mitrevsk2021ontology}, scene graphs~\cite{johnson2015image,armeni20193d}, and other relational structures~\cite{brachman2004knowledge,russell1995modern}. Such representations are useful not only because they are interpretable, but because individual facts can change which robot actions are feasible. In STRIPS-, PDDL-, and FDR-style planning, for example, predicates describe the state and action preconditions and effects define its transition structure~\cite{fikes1971strips,aeronautiques1998pddl,helmert2009concise}. Classical planners can then compute executable action sequences over this structured model~\cite{hoffmann2001ff,helmert2006fast,correa2024lifted}. This connection between state facts and action feasibility is central to our setting: uncertainty about whether an object is \texttt{plastic} or a tomato is \texttt{ripe} can alter the subsequent sorting or harvesting plan.

Standard symbolic planning generally assigns each fact a single truth value and therefore does not directly represent noisy observations, hidden states, or competing interpretations of the world. Probabilistic logical models address this limitation by attaching uncertainty to facts or rules~\cite{richardson2006markov,qu2019probabilistic}, while commonsense and relational inference can infer unobserved facts from spatial, semantic, or contextual knowledge~\cite{tenorth2012knowledge,kunze2014using,hasegawa2023inferring,daruna2019robocse,mota2021integrated,zhang2024icorpp}. These approaches enrich the state representation, but inference alone does not determine how unresolved uncertainty should influence sequential action selection or when it should be resolved through external information.

POMDPs provide a standard decision-theoretic account of this problem by maintaining a belief over possible states and planning with respect to future actions and observations~\cite{kaelbling1998planning,pineau2003point,ross2008online,somani2013despot}. Exact belief-space planning is difficult in large domains, motivating point-based, tree-search, and sampling-based approximations~\cite{pineau2003point,ross2008online,silver2010monte,bai2015intention}. Particle methods represent the belief using sampled complete states and have been applied to navigation, exploration, object search, manipulation, autonomous driving, and aerial robotics~\cite{lauri2016exploration,kim2021plgrim,xiao2019objectsearch,pajarinen2017multiobject,pajarinen2022composition,brechtel2014continuous,bai2011uav,vanegas2016uav}. Factored approximations instead exploit structure among state variables to reduce inference cost~\cite{hansen2000dynamic,feng2001approximate,veiga2014point,murphy2013factored}. For symbolic tasks whose predicates are coupled through preconditions, effects, and mutual constraints, complete-state hypotheses offer a direct way to preserve planning-relevant dependencies while still supporting sampling-based online planning.

Several approaches combine symbolic structure with probabilistic planning. CORPP constructs priors over possible worlds from commonsense knowledge and integrates them with a POMDP planner~\cite{zhang2015corpp,kurniawati2008sarsop}; other methods exploit symbolic structure, abstraction, or hierarchical state spaces to manage planning complexity~\cite{sanner2010symbolic,serrano2021knowledge,li2026uncertainty}. More recent systems use language models to provide open-ended task knowledge, generate state hypotheses, or produce formal specifications that can be executed by downstream planners~\cite{tang2025tru,huang2023voxposer,dai2024optimal,liu2025delta,paulius2025bootstrapping,quartey2025verifiably}. Related work also grounds language-model-generated plans through affordances, programmatic interfaces, scene representations, retrieval, and execution feedback~\cite{ahn2022can,huang2022inner,singh2023progprompt,liang2023code,ding2023task,gu2024conceptgraphs,leanza2025conceptbot,singh2025adaptbot,xu2024p,wang2025instructrag}. These systems broaden the knowledge available to a robot, whereas our focus is narrower: given competing symbolic state hypotheses and an explicit transition model, we study how the robot should acquire external information that resolves planning-relevant ambiguity.

\subsection{Selective Information Acquisition under Uncertainty}~\label{rel:selective_information}
A robot can reduce uncertainty by gathering additional observations, changing its viewpoint, exploring the environment, or consulting an external source. Active perception and information-gathering planners select actions according to their expected effect on the belief or future task utility. Examples include exploration based on information gain, online POMDP planning for object search, and belief-aware manipulation~\cite{lauri2016exploration,kim2021plgrim,xiao2019objectsearch,pajarinen2022composition}. These methods usually treat information acquisition as a robot action. Our setting considers a complementary channel: some task facts are difficult or costly for the robot to infer reliably, but can be verified through targeted external feedback.

A simple alternative is continuous supervision or human-initiated correction~\cite{kelly2019hg,mandlekar2020human}. This can provide timely recovery when an operator detects a failure, but it requires the operator to monitor execution and decide when intervention is necessary. Sustained monitoring can increase workload and is vulnerable to vigilance decrement and inattentional blindness~\cite{davies1982psychology,davies2013human,mack2003inattentional,wong2017workload,galster2006managing}. System-initiated interaction moves the detection decision to the robot: rather than waiting for a person to notice a problem, the robot requests assistance when its internal estimate indicates that autonomous execution may be unreliable.

Uncertainty-aware help-seeking systems implement this idea using policy confidence, task difficulty, trust, or calibrated prediction sets~\cite{singi2023hula,igbinedion2024learning,zhao2025conformalized,mangalindan2025seek}. Embodied agents and collaborative systems likewise ask clarification questions when model confidence is insufficient~\cite{shen2024elba,testoni2024asking}. This line of work establishes that \emph{when} to interact can be treated as a decision rather than a fixed supervision schedule. It also exposes a trade-off: querying too rarely leaves consequential uncertainty unresolved, whereas querying too often adds interaction without necessarily improving task performance. Our timing decision follows the same high-level principle but applies it to a belief over symbolic world states, allowing the query trigger to reflect ambiguity in facts that constrain feasible plans.

\subsection{Action-Level Clarification and State Verification}~\label{rel:query_content}
The closest related approaches combine uncertainty-aware timing with action- or plan-level clarification. KnowNo uses conformal prediction to construct a set of candidate robot actions and asks the user to select among them when the prediction set is not sufficiently resolved~\cite{ren2023robots}. Related language-model planning systems use calibrated plan selection, retrieval-augmented reasoning, or uncertainty over candidate plans to decide when help is required~\cite{liang2024introspective,mullen2025lbap}. These approaches provide a principled interface between uncertainty estimation and system-initiated interaction, and action guidance can be highly efficient when the correct action or plan appears among the candidates.

Our work studies a different query target. Action- and plan-level queries ask the external source to participate directly in selecting what the robot should do next. State-level verification instead asks whether a planning-relevant fact is true and uses the answer as evidence for updating the robot's belief. The planner then determines the next action from the refined state estimate. The distinction matters when action candidates are generated from an incorrect or incomplete state estimate: additional guidance over the resulting candidates may not resolve the upstream state ambiguity on which later decisions also depend. Conversely, state verification may require more than one question before the belief is sufficiently resolved. The two interfaces therefore present different trade-offs rather than a universal dominance relation.

Query content introduces a second decision beyond query timing. Asking about every uncertain fact would reproduce dense supervision at the level of state variables, while an arbitrary question may provide little information about the competing hypotheses. Active query selection instead prioritizes information according to its expected effect on the current belief or downstream decision. In our framework, expected posterior entropy determines which symbolic predicate to verify, and the response filters inconsistent world hypotheses. This connects the content of feedback to the same belief used for planning and permits the effects of \emph{when} and \emph{what} to ask to be evaluated separately.

The form of a query also changes what the user must understand and decide. Action selection can require interpreting the current situation and anticipating the consequences of candidate actions, whereas fact verification asks for evidence about a stated property of the world. Prior research shows that workload and situation awareness can affect intervention quality in automated systems~\cite{endsley2018automation,murata2000ergonomics,borragan2017cognitive,rajavenkatanarayanan2019towards}. We therefore consider user-facing effort alongside task success and query efficiency, without treating human supervision itself as the primary object of study.

In summary, prior work provides strong foundations for symbolic task representation, belief-space planning, and uncertainty-triggered assistance. The remaining question addressed here is how to use a shared belief over symbolic states to coordinate both feedback decisions: request information when unresolved state ambiguity threatens reliable execution, and select the state fact expected to resolve that ambiguity efficiently. This formulation keeps external feedback focused on evidence about the world while leaving subsequent action selection to the robot.
